\chapter{VR-Forms}
\label{chap:forms}
\uses{chap:setsop, chap:sets, chap:setszfa}

\section{Position}
\label{sec:forms_position}

VR-Forms is the two-register apparatus over the operational sets.  The
\emph{operational register} (L\textsubscript{0}) describes operational sets
in the sense of Chapter~\ref{chap:setsop}: revealing functionalities,
identical when bisimilar by witness.  The \emph{formal register}
(L\textsubscript{1}) accommodates syntactic descriptions without operational
obligation---classically uncountable objects, paradoxical classes, and any
correctly built term regardless of operationality or consistency.  The two
registers are connected by the translation $\pi$ and by Theorem~III.1
(conservativity).

\textbf{Clarification on register language.} The registers describe modes
of description, not separate operational levels.  All descriptions are
operational acts; the distinction is whether the described referent has an
operational correlate.

\textbf{Two readings of the operational register (2026-09-12, integrity
programme, step~3).}  Until 2026-09-12 the operational register of the Lean
formalisation was Mathlib's well-founded ZFC universe \texttt{OSet = ZFSet}
(Chapter~\ref{chap:sets}), and every object above the language layer
carried \texttt{[propext, Quot.sound]}.  The formalisation now lives in the
core library \texttt{VR} over the operational universe \texttt{OpSet}
(Chapter~\ref{chap:setsop}) --- \texttt{VR/Forms/Realisability.lean},
\texttt{Transit.lean}, \texttt{Examples.lean}, \texttt{Substrate.lean} ---
and every declaration is on the empty axiom list, enforced by the core's
build-time guard.  The ZFC reading is kept, named as such, in
\texttt{VRClassical/Forms/Bridge.lean}: the predicate
\texttt{isRealisableZFC} (\S\,\ref{sec:forms_bridge}) is realisability with
the operational register taken to be \texttt{ZFSet}.  The two readings
disagree on exactly one named term, $\ulcorner\mathrm{AFA}\urcorner$, and
their disagreement is itself a theorem (\S\,\ref{sec:forms_examples}): the
ZFA boundary of VR-Sets is a boundary of \emph{groundedness} (a mode), not
of operationality.  Axioms that appear in the bridge are Mathlib's
(\texttt{ZFSet}/\texttt{PSet}) --- the bridge's, not VR's.

\textbf{Architectural note.} The Lean formalisation is a \emph{shallow
embedding}: formal terms are Lean strings carrying a register tag; $\pi$ is a
total Lean function; realisability is a Lean predicate.  Theorem~III.1
(conservativity), proved in the preprint, is formalised beside the shallow
embedding by a separate deep embedding
(\texttt{Forms/Conservativity.lean}, \texttt{ConservativityFOL.lean},
\texttt{ConservativityComprehension.lean}): \texttt{Formula} over de Bruijn
terms, a classical Hilbert \texttt{Provable}, the translation $\pi$, and the
conservativity theorem by induction on derivations, with a concrete VR
instance.  See \S\,\ref{subsec:conservativity}.

\noindent\axiomprofile{[]} for every object of the core files (language,
realisability, transit, conservativity in all three storeys, examples,
substrate); the ZFC bridge inherits Mathlib's
\axiomprofile{[propext, Quot.sound]} and, through \texttt{PSet}, no
\texttt{Classical.choice}.

%──────────────────────────────────────────────────────────────────────
\section{Formal language (Part II)}
\label{sec:forms_language}

\begin{definition}[Two registers]
\label{def:Register}
\lean{VR.Forms.Register}
\leanok
\texttt{Register} is an inductive type with two constructors:
\texttt{.operational} (the operational register~L\textsubscript{0}) and
\texttt{.formal} (the formal register~L\textsubscript{1}).
\axiomprofile{[]} --- no imports; pure Lean~4 prelude.
\end{definition}

\begin{definition}[Formal terms]
\label{def:FormalTerm}
\lean{VR.Forms.FormalTerm}
\leanok
\uses{def:Register}
\texttt{FormalTerm} is a structure with fields \texttt{description :
String} and \texttt{register : Register}.  The notation
$\ulcorner\tau\urcorner$ always produces
\texttt{FormalTerm.mk\,$\tau$\,.formal}.  The Principle of Forms (§\,II.4):
any syntactically correct string specifies a formal term --- no
operationality, consistency, or constructivity is required.
Injectivity of the constructor is proved by hand
(\texttt{FormalTerm.mk\_eq\_iff}); the auto-generated \texttt{injEq} is
not produced, since it would carry \texttt{propext}.
\axiomprofile{[]} --- Language.lean imports nothing.
\end{definition}

%──────────────────────────────────────────────────────────────────────
\section{Realisability (Part II \S\,II.7)}
\label{sec:forms_realisable}

\begin{definition}[Operational realisability predicate]
\label{def:isRealisable}
\lean{VR.Forms.isRealisable, VR.Forms.isRealisableN, VR.Forms.named}
\leanok
\uses{def:FormalTerm, def:setsop_opset}
Realisability is defined in two steps.  A closed-world classifier
$\texttt{named} : \texttt{FormalTerm} \to \texttt{Named}$ sends a term, by
decidable equality of its description, to one of
$\{\texttt{empty}, \texttt{omega}, \texttt{pair}, \texttt{afa},
\texttt{other}\}$; then
$\texttt{isRealisableN} : \texttt{Named} \to \mathrm{Prop}$ gives each class
its existential proposition over \texttt{OpSet}:

\begin{tabular}{ll}
$\ulcorner$\texttt{"∅"}$\urcorner$ & $\exists\, s : \mathrm{OpSet},\; \forall x,\; \lnot\, x \mathrel{\mathtt{Mem}} s$ \\
$\ulcorner$\texttt{"omega"}$\urcorner$ & $\exists\, s : \mathrm{OpSet},\; \mathtt{vn}\,0 \mathrel{\mathtt{Mem}} s \land \forall n,\; n \mathrel{\mathtt{Mem}} s \to \mathtt{succ}\,n \mathrel{\mathtt{Mem}} s$ \\
$\ulcorner$\texttt{"pair"}$\urcorner$ & $\forall a\,b : \mathrm{OpSet},\; \exists\, s,\; \forall x,\; x \mathrel{\mathtt{Mem}} s \leftrightarrow x \mathrel{\mathtt{Equiv}} a \lor x \mathrel{\mathtt{Equiv}} b$ \\
$\ulcorner$\texttt{"AFA"}$\urcorner$ & $\forall\, (V : \mathrm{Type})\,(E : V\to V\to\mathrm{Prop}),\; \exists\, d : V \to \mathrm{OpSet},\; \mathtt{IsDecoration}\,E\,d$ \\
any other & \texttt{False}
\end{tabular}

$\texttt{isRealisable}\,t := \texttt{isRealisableN}\,(\texttt{named}\,t)$.
Membership is read as \texttt{Mem} and identity as \texttt{Equiv}: the pair
$\{a,b\}$ has members \emph{identical to} $a$ or $b$ by witness, not by
extension.  Each realisable case is an existential Lean Prop (not
\texttt{True}), preserving the mathematical content of the witness.  The
two-step form replaces the earlier \texttt{match} on string literals: a
\texttt{match} with overlapping literal patterns generates a splitter that
carries \texttt{propext}; the \texttt{if}-chain of \texttt{named} does not.
\axiomprofile{[]}.
\end{definition}

\begin{theorem}[Base realisability lemmas]
\label{thm:isRealisable_base}
\lean{VR.Forms.isRealisable_empty, VR.Forms.isRealisable_omega,
      VR.Forms.isRealisable_pair, VR.Forms.isRealisable_AFA}
\leanok
\uses{def:isRealisable, thm:setsop_basic, thm:setsop_omega, thm:setsop_afa}
$\texttt{isRealisable\_empty}$: witness \texttt{OpSet.empty}, which reveals
nothing (\texttt{not\_mem\_empty}).
$\texttt{isRealisable\_omega}$: witness \texttt{OpSet.omega}
(\texttt{empty\_mem\_omega}, \texttt{omega\_succ\_closed}).
$\texttt{isRealisable\_pair}$: witness \texttt{OpSet.pair a b}
(\texttt{mem\_pair}, up to \texttt{Equiv}).
$\texttt{isRealisable\_AFA}$: witness \texttt{OpSet.decorate E}
(\texttt{decorate\_isDecoration}) --- every graph has a decoration in the
operational universe.  All four on \axiomprofile{[]}.
\end{theorem}

\begin{remark}[Three-category structure of formal terms]
\label{rem:three_category}
\lean{VR.Forms.isRealisable, VR.Forms.isRealisableZFC}
\leanok
\uses{def:isRealisable, def:Conjecture_IV_1, def:Conjecture_IV_2,
      def:AFA_Statement}
Across the two readings the formal terms fall into three categories:
(a)~\emph{Provably realisable} in the operational universe ---
$\ulcorner$\texttt{"∅"}$\urcorner$, $\ulcorner$\texttt{"omega"}$\urcorner$,
$\ulcorner$\texttt{"pair"}$\urcorner$, $\ulcorner$\texttt{"AFA"}$\urcorner$,
with concrete \texttt{OpSet} witnesses;
(b)~\emph{Open realisability} in the ZFC reading ---
$\ulcorner$\texttt{"Conjecture\_IV\_X\_Statement"}$\urcorner$: realisability
is mathematically open, mirroring VR-Sets Conjectures~IV.1/IV.2 (only
\texttt{isRealisableZFC} names these terms; in the operational reading they
fall to the catch-all);
(c)~\emph{Provably non-realisable} --- split into trivially-False (the
catch-all: Russell, Vitali, classical~$\mathbb{R}$, classical
$\mathcal{P}(\mathbb{N})$, no set theorem needed) and ZFC-refutable
($\ulcorner$\texttt{"AFA\_Statement"}$\urcorner$ via
\texttt{AFA\_Refuted}, in the ZFC reading only).  The term
$\ulcorner\mathrm{AFA}\urcorner$ is the one that changes category with the
reading --- see Theorem~\ref{thm:mixed_AFA_two_registers}.
\end{remark}

%──────────────────────────────────────────────────────────────────────
\section{Transit (Part III)}
\label{sec:forms_transit}

\begin{definition}[Translation $\pi$ (shallow)]
\label{def:translate_pi}
\lean{VR.Forms.translate_pi, VR.Forms.translateN}
\leanok
\uses{def:FormalTerm, def:isRealisable, thm:setsop_basic}
$\texttt{translate\_pi} : \texttt{FormalTerm} \to \mathrm{Prop}$ is the
\emph{specific} layer (rule~(ii) of §\,III.2): for each named formal term,
the defining predicate of the concrete operational set.  For
$\ulcorner$\texttt{"∅"}$\urcorner$: $\forall x,\; \lnot\, x
\mathrel{\mathtt{Mem}} \texttt{OpSet.empty}$ (not an existential).  For
$\ulcorner$\texttt{"omega"}$\urcorner$: $\mathtt{vn}\,0 \mathrel{\mathtt{Mem}}
\texttt{OpSet.omega} \land \ldots$.  For $\ulcorner$\texttt{"pair"}$\urcorner$:
$\forall a\,b\,x,\; x \mathrel{\mathtt{Mem}} \texttt{OpSet.pair}\,a\,b
\leftrightarrow x \mathrel{\mathtt{Equiv}} a \lor x \mathrel{\mathtt{Equiv}}
b$.  For $\ulcorner$\texttt{"AFA"}$\urcorner$: $\forall V\,E,\;
\mathtt{IsDecoration}\,E\,(\texttt{OpSet.decorate}\,E)$.  Catch-all:
\texttt{False}.  Contrast with \texttt{isRealisable} (\emph{existential}
layer).  \axiomprofile{[]}.
\end{definition}

\begin{theorem}[π-translation lemmas]
\label{thm:translate_pi_base}
\lean{VR.Forms.translate_pi_empty, VR.Forms.translate_pi_omega,
      VR.Forms.translate_pi_pair, VR.Forms.translate_pi_AFA}
\leanok
\uses{def:translate_pi, thm:setsop_basic, thm:setsop_omega, thm:setsop_afa}
Direct proofs for the four named terms by the SetsOp theorems
\texttt{not\_mem\_empty}, \texttt{empty\_mem\_omega} with
\texttt{omega\_succ\_closed}, \texttt{mem\_pair},
\texttt{decorate\_isDecoration}.  All on \axiomprofile{[]}.
\end{theorem}

\begin{theorem}[Transit: $\pi$ implies realisability]
\label{thm:translate_implies_realisable}
\lean{VR.Forms.translate_implies_realisable, VR.Forms.translateN_implies_realisableN}
\leanok
\uses{def:translate_pi, def:isRealisable}
$\forall t : \texttt{FormalTerm},\;
\texttt{translate\_pi}(t) \Rightarrow \texttt{isRealisable}(t)$.
Proved class by class on \texttt{Named}
(\texttt{translateN\_implies\_realisableN}): in each realisable class the
specific predicate names the witness, and existential introduction closes
the goal; the catch-all is vacuous.  The \emph{converse does not hold}:
from $\exists\, s,\; \ldots$ one cannot recover the specific named object
without Skolemisation --- the T$\to$O asymmetry (TR-R1).  The transit
pattern operates exclusively in the forward direction.  \axiomprofile{[]}.
\end{theorem}

%──────────────────────────────────────────────────────────────────────
\subsection{The conservativity boundary (Theorem~III.1)}
\label{subsec:conservativity}

\begin{theorem}[Conservativity of $T_1$ over $T_0$]
\label{thm:conservativity}
\lean{VR.Forms.ConservativityFOL.conservativity}
\leanok
\uses{def:translate_pi, def:isRealisable}
\textbf{Theorem III.1} (preprint §\,III.2, verbatim): «The theory $T_1$
(VR-Forms) is conservative over $T_0$ (VR-Sets) in the operational
register: any formula $\varphi \in L_0$ that is derivable in $T_1$ is
already derivable in $T_0$.  Equivalently: formal terms do not produce
new operational theorems.»

\textbf{Formalised} (2026-06) by a deep embedding via relative interpretation, in
\texttt{Forms/Conservativity.lean} (propositional floor) and
\texttt{Forms/ConservativityFOL.lean} (full FOL).  The two-register language $L_1$ consists of
operational atoms and $0$-ary formal atoms over de Bruijn terms (variables, constants such as
$\emptyset$, and $n$-ary function symbols such as $t$); provability is classical Hilbert
(K, S, Peirce, MP, $\forall$-elim, $\forall$-distribution, generalization); the translation
$\pi$ sends each formal atom to its closed operational meaning.  The proof is by induction on
derivations --- $\pi$ commutes with every rule, and \texttt{piTr\_subst} shows $\pi$ commutes
with substitution (kept trivial by $0$-ary formal atoms and closed $\pi$-images, so no de-Bruijn
substitution composition arises).  A concrete VR instance (\texttt{VRExample}: $\emptyset$,
$\in$, $t$, the formal term $\ulcorner\emptyset\urcorner \mapsto \forall x\,\lnot(x\in\emptyset)$)
discharges the hypothesis end-to-end (\texttt{conservativity\_empty\_concrete}).

\textbf{Comprehension storey (full fidelity to the preprint's $\pi$, 2026-06).}  The FOL floor
models formal atoms as $0$-ary; the preprint's $\pi$ (\S\,III.2, rule for $\ulcorner x\in\tau\urcorner$
with $\tau=\{y:\psi(y)\}$) is the genuine comprehension translation $\pi(\ulcorner x\in\tau\urcorner)=\psi(x)$,
which substitutes the argument.  This is machine-checked in
\texttt{Forms/ConservativityComprehension.lean}: terms and formulas are \emph{mutually} recursive
($\mathtt{Tm}$ carries set-builders $\{y:\psi\}$), so $\pi$ substituting $\psi(x)$ forces the full
de-Bruijn substitution calculus --- $\mathtt{subst\_lift}$, $\mathtt{lift\_lift}$,
$\mathtt{lift\_subst}$, and the substitution lemma $\mathtt{subst\_subst}$.
The keystone $\mathtt{piFml\_subst}$ ($\pi$ commutes with substituting a variable) closes precisely
via $\mathtt{subst\_subst}$ on the comprehension case; conservativity then follows as before, with a
non-vacuous instance whose hypothesis is proved ($\pi$ of a comprehension axiom is the Hilbert theorem
$X\to X$).  This is \emph{course A}: set-builders occur only as the right member of $\in$, so the left
member of every membership is operational (a variable).  A \emph{total} $\pi$ unfolding set-builders in
every position (course B) is provably ill-defined --- it diverges on self-membered formal terms (the
Russell term $\ulcorner\{x:x\notin x\}\urcorner$ gives $\pi(\ulcorner R\in R\urcorner)=\lnot\pi(\ulcorner R\in R\urcorner)$,
no fixpoint).  This is the formal-register counterpart of why VR holds paradoxes safely: they cannot be
unfolded into operational theorems --- $\pi$ is simply undefined on them, not defined-and-false.

\textbf{Axiom profile (2026-09-12 sweep).}  All three storeys are on \axiomprofile{[]}: the de Bruijn
lemmas that had been closed by \texttt{simp}/\texttt{omega} (both reach \texttt{propext}) were
re-proved by hand (\texttt{NatAux}, explicit \texttt{rw [if\_pos …]}), and the auto-generated
\texttt{injEq} lemmas are switched off.

\begin{remark}[Scope and remaining polish]
\label{rem:conservativity_boundary}
What was the cycle's one premeditated boundary --- a proof-theory project beyond the shallow
embedding --- is a machine-checked result.  Two documented refinements remain, by choice,
and are polish rather than gaps in the theorem: \texttt{gen} is stated without the eigenvariable
side-condition (it would match standard first-order soundness, but is not needed for the
syntactic $\pi$-transport that conservativity is), and the worked VR instance covers $\emptyset$
(richer instances such as $\omega$ are marginal).
\end{remark}
\end{theorem}

%──────────────────────────────────────────────────────────────────────
\section{The ZFC reading and the bridge theorems (Part V)}
\label{sec:forms_bridge}

The bridge file \texttt{VRClassical/Forms/Bridge.lean} keeps the reading
VR-Forms had until 2026-09-12: realisability with the operational register
taken to be Mathlib's \texttt{OSet = ZFSet}.  It is a second predicate,
\texttt{isRealisableZFC}, over the same formal terms; its three positive
cases are the ZFC-era witnesses (\texttt{osetEmpty}, \texttt{omega\_OSet},
\texttt{osetPair}, Chapter~\ref{chap:sets}), and it alone names the three
statements about Mathlib's universes.

\begin{theorem}[Negative bridge: AFA is non-realisable in the ZFC reading]
\label{thm:bridge_AFA}
\lean{VR.Forms.bridge_AFA, VR.Forms.isRealisableZFC}
\leanok
\uses{def:isRealisable, thm:AFA_Refuted}
$\lnot\,\texttt{isRealisableZFC}\,\ulcorner$\texttt{"AFA\_Statement"}$\urcorner$.
Proof: \texttt{isRealisableZFC} reduces definitionally to
\texttt{VR.Sets.AFA\_Statement}; the theorem is exactly
\texttt{AFA\_Refuted} from Chapter~\ref{chap:sets}.  The cross-cycle chain:
\texttt{PSet.mem\_irrefl} $\to$ \texttt{AFA\_Refuted} $\to$
\texttt{bridge\_AFA}.  This is a statement about Mathlib's well-founded
universe; in the operational universe the same description is realised
(Theorem~\ref{thm:isRealisable_base}, \texttt{isRealisable\_AFA}).
\axiomprofile{[propext, Quot.sound]} (Mathlib's, via \texttt{ZFSet}).
\end{theorem}

\begin{theorem}[Open conditional bridges: Conjectures~IV.1 and IV.2]
\label{thm:bridge_Conjectures}
\lean{VR.Forms.bridge_Conjecture_IV_1, VR.Forms.bridge_Conjecture_IV_2}
\leanok
\uses{def:isRealisable, def:Conjecture_IV_1, def:Conjecture_IV_2}
$\texttt{isRealisableZFC}\,\ulcorner$\texttt{"Conjecture\_IV\_1\_Statement"}$\urcorner
\leftrightarrow \texttt{Conjecture\_IV\_1\_Statement}$, and likewise for
IV.2.  Both iff proofs are $\langle\texttt{id},\texttt{id}\rangle$
(definitional equality); the mathematical \emph{content} is open.  Status:
open realisability (category~(b) of the triadic structure), in the ZFC
reading.
\end{theorem}

%──────────────────────────────────────────────────────────────────────
\section{Examples and mixed formulas (Parts V--VII)}
\label{sec:forms_examples}

\begin{theorem}[Four non-realisable formal terms]
\label{thm:not_realisable}
\lean{VR.Forms.not_isRealisable_Russell, VR.Forms.not_isRealisable_Vitali,
      VR.Forms.not_isRealisable_classical_R,
      VR.Forms.not_isRealisable_classical_powerset_N}
\leanok
\uses{def:isRealisable}
$\lnot\,\texttt{isRealisable}\,\ulcorner$\texttt{"Russell\_class"}$\urcorner$,
$\lnot\,\texttt{isRealisable}\,\ulcorner$\texttt{"Vitali"}$\urcorner$,
$\lnot\,\texttt{isRealisable}\,\ulcorner$\texttt{"classical\_R"}$\urcorner$,
$\lnot\,\texttt{isRealisable}\,\ulcorner$\texttt{"classical\_powerset\_N"}$\urcorner$.
All four proofs are \texttt{id}: the closed world of \texttt{named} sends
each description to \texttt{other}, hence to \texttt{False}.
Non-realisability is metatheoretic (category~(c), Level~1): no set theorem
is needed.  Contrast with \texttt{bridge\_AFA} (Level~2: refutable in the
ZFC reading).  \axiomprofile{[]}.
\end{theorem}

\begin{theorem}[Positive mixed formula: $\omega$ in both registers]
\label{thm:mixed_omega}
\lean{VR.Forms.mixed_omega_two_register}
\leanok
\uses{def:isRealisable, def:translate_pi, thm:translate_pi_base}
$\mathtt{vn}\,0 \mathrel{\mathtt{Mem}} \texttt{OpSet.omega} \land
\texttt{isRealisable}\,\ulcorner$\texttt{"omega"}$\urcorner$.
Proof: $\langle\texttt{OpSet.empty\_mem\_omega},\,
\texttt{isRealisable\_omega}\rangle$.
Both registers affirm $\omega$: operational membership and formal
realisability agree on the same object.  Mixed formulas combine operational
statements about \texttt{OpSet} with realisability predicates on formal
terms --- Lean \texttt{Prop}s at the meta level, not a third register.
\axiomprofile{[]}.
\end{theorem}

\begin{theorem}[Central junction: the ZFA boundary is a boundary of mode]
\label{thm:mixed_AFA_two_registers}
\lean{VR.Forms.mixed_AFA_two_registers, VR.Forms.afa_two_registers}
\leanok
\uses{thm:isRealisable_base, thm:bridge_AFA, thm:setsop_afa, thm:setsop_grounded}
$\texttt{isRealisable}\,\ulcorner$\texttt{"AFA"}$\urcorner \land
\lnot\,\texttt{isRealisableZFC}\,\ulcorner$\texttt{"AFA\_Statement"}$\urcorner$.
Proof: $\langle\texttt{isRealisable\_AFA},\,\texttt{bridge\_AFA}\rangle$
(\texttt{afa\_two\_registers}).

\emph{Junction theorem of VR-Forms.}  The formal register's anti-foundation
description has an operational correlate --- \texttt{OpSet.afa}: every
graph is decorated in the operational universe, since a set \emph{is} a
pointed graph and identity is bisimulation --- while its ZFC-register
reading is refuted in Mathlib's well-founded universe
(\texttt{PSet.mem\_irrefl} $\to$ \texttt{AFA\_Refuted}).  Both conjuncts are
theorems; they speak of different universes.  What Chapter~\ref{chap:sets}
recorded as Boundary~B.5 (``ZFA total absence'') is thereby located: it is
a boundary of \emph{groundedness} --- \texttt{IsGrounded} is a predicate on
the graph, the ZFC-mode (\S\,\ref{sec:setsop_grounded}) --- not a boundary
of operationality.  See Observation~9.
\end{theorem}

%──────────────────────────────────────────────────────────────────────
\section{Substrate (Stage 6)}
\label{sec:forms_substrate}

\begin{definition}[Carrier and operational substrate]
\label{def:Carrier}
\lean{VR.Forms.Carrier, VR.Forms.Operational}
\leanok
\uses{def:setsop_opset, thm:setsop_equiv, def:FormalTerm}
$\texttt{Carrier}$ is an inductive type with two constructors:
$\texttt{.obj} : \mathrm{OpSet}.\{0\} \to \texttt{Carrier}$ (operational
objects in the operational register) and
$\texttt{.term} : \texttt{FormalTerm} \to \texttt{Carrier}$ (formal terms).

$\texttt{Operational} : \texttt{Carrier} \to \mathrm{Prop}$ by:
$\texttt{Operational}(\texttt{.obj}\,a) := a \mathrel{\mathtt{Equiv}} a$
--- the substrate of an operational set is the act it \emph{is}, its
revealing functionality, witnessed minimally by the identity bisimulation
(the diagonal relation with its two simulation proofs: a construction, not
\texttt{True});
$\texttt{Operational}(\texttt{.term}\,\_) := \texttt{True}$ --- for a formal
term the substrate is the act of inscription (Principle of Forms, §\,II.4).

\textbf{Why not well-foundedness.}  The 2026-05-28 gate chose
$\texttt{Acc}({\cdot}\in{\cdot})$ as the substrate predicate on
\texttt{OSet = ZFSet}, where it is total because Mathlib's universe is
well-founded by construction.  On \texttt{OpSet} that choice is
\emph{refuted} as a substrate: AFA holds there, so groundedness
(\texttt{IsGrounded}) is not total --- it is a mode, separately available.
The frozen statement («operationality is a total substrate; every operand
and every result, including the empty result $\emptyset$, carries
operational substrate») survives with the correct predicate.
\end{definition}

\begin{theorem}[Substrate totality and base cases]
\label{thm:operational_total}
\lean{VR.Forms.operational_total, VR.Forms.operational_empty,
      VR.Forms.omega_substrate}
\leanok
\uses{def:Carrier, thm:setsop_equiv, thm:setsop_omega}
$\texttt{operational\_total}$: $\forall c : \texttt{Carrier},\;
\texttt{Operational}(c)$.  Proof: \texttt{.obj}~branch by
\texttt{OpSet.Equiv.refl}; \texttt{.term}~branch by \texttt{trivial}.

$\texttt{operational\_empty}$: $\texttt{Operational}(\texttt{.obj}\,
\texttt{OpSet.empty})$ --- the primordial object $\emptyset$, the limit case
the curator insisted on, carries operational substrate.

$\texttt{omega\_substrate}$: $\texttt{Operational}(\texttt{.obj}\,
\texttt{OpSet.omega}) \land \texttt{Operational}(\texttt{.term}\,
\ulcorner$\texttt{"omega"}$\urcorner)$ --- $\omega$ carries substrate in
both registers, connecting Stage~6 to
\texttt{mixed\_omega\_two\_register}.

\noindent\axiomprofile{[]} for all three.
\end{theorem}

%──────────────────────────────────────────────────────────────────────
\section{Conservativity: from boundary to formalised result (\S\,IX.2)}
\label{sec:forms_boundary}

VR-Forms Lean originally had \emph{one} explicit structural boundary (at conservativity),
in contrast to VR-Sets Lean's five mathematical-content boundaries.  That boundary has
since been \textbf{closed}: conservativity is formalised by a deep embedding
(\S\,\ref{subsec:conservativity}).

\begin{remark}[Observation~3 --- Conservativity, formalised]
\label{rem:obs3_boundary}
\lean{VR.Forms.ConservativityFOL.conservativity}
\leanok
\uses{thm:conservativity}
Theorem~III.1 was mathematically proved in the preprint (full inductive proof via $\pi$);
the shallow embedding deliberately did not formalise it, documenting the boundary in
\texttt{Transit.lean} and declining a \texttt{def Conjecture\_Conservativity} (which would
misrepresent a proved result as open).  It is machine-checked by a separate deep
embedding (\texttt{Forms/ConservativityFOL.lean}): a classical Hilbert calculus over de
Bruijn terms with $n$-ary functions and predicates, the $\pi$-translation, and conservativity
by induction on derivations, with a concrete VR $\emptyset$ instance --- on the empty axiom
list.  The shallow embedding remains the working layer; the deep embedding stands beside it
as the proof.
\end{remark}

\begin{remark}[Observation~4 --- Closed-world classification instead of a literal \texttt{match}]
\label{rem:obs4_catchall}
\lean{VR.Forms.named, VR.Forms.translate_implies_realisable}
\leanok
\uses{thm:translate_implies_realisable}
The original \texttt{match} on description strings had two costs: its
catch-all \texttt{|\,\_ => False} did not reduce in term mode for a
schematic variable, forcing \texttt{by\_cases} discrimination; and the
splitter Lean generates for overlapping literal patterns carries
\texttt{propext} (measured in the 2026-09-12 sweep).  The classifier
\texttt{named} (an \texttt{if}-chain over \texttt{DecidableEq FormalTerm}
into the finite \texttt{Named}) removes both: every theorem about
\texttt{isRealisable} is a case analysis on five constructors, and the
module is on \texttt{[]}.
\end{remark}

\begin{remark}[Observation~5 --- Realisability and $\pi$-translation are complementary layers]
\label{rem:obs5_layers}
\lean{VR.Forms.isRealisable, VR.Forms.translate_pi,
      VR.Forms.translate_implies_realisable}
\leanok
\uses{def:isRealisable, def:translate_pi, thm:translate_implies_realisable}
\texttt{isRealisable} (existential: $\exists$ witness) and
\texttt{translate\_pi} (specific: named concrete object) form two
complementary layers connected by \texttt{translate\_implies\_realisable}
via existential introduction.  The converse (specific $\Leftarrow$
existential) is not provable: Skolemisation is unavailable from the
existential in the shallow embedding.  The transit pattern of §\,IV.2
operates exclusively in the forward direction --- the same asymmetry that
governs the apparatus (Chapter~\ref{chap:apparatus}, T$\to$O has no
mechanism).
\end{remark}

\begin{remark}[Observation~6 --- Two-level structure of negative cases]
\label{rem:obs6_negative}
\lean{VR.Forms.not_isRealisable_Russell, VR.Forms.bridge_AFA}
\leanok
\uses{thm:not_realisable, thm:bridge_AFA}
Two levels of non-realisability:
\emph{Level~1 (trivially False)}: Russell, Vitali, classical~$\mathbb{R}$,
classical~$\mathcal{P}(\mathbb{N})$ --- the closed world sends them to
\texttt{False}, proof is \texttt{id}, no set theorem needed; this level lives
in the core on \texttt{[]}.
\emph{Level~2 (refutable in a reading)}: $\ulcorner$\texttt{AFA\_Statement}$\urcorner$
--- \texttt{isRealisableZFC} returns a mathematically formulated
\texttt{Prop}; non-realisability requires \texttt{AFA\_Refuted} from
VR-Sets, and lives in the bridge.  The operational reading has no Level~2:
every description it names is realised.
\end{remark}

\begin{remark}[Observation~7 --- Three-category structure mirrors VR-Sets tier structure]
\label{rem:obs7_triad}
\lean{VR.Forms.bridge_AFA, VR.Forms.bridge_Conjecture_IV_1,
      VR.Forms.isRealisable_empty}
\leanok
\uses{thm:isRealisable_base, thm:bridge_AFA, thm:bridge_Conjectures}
The triadic classification of formal terms (provably realisable~/ open ~/
provably non-realisable) parallels VR-Sets Stage~11's three-tier
formalisation result (proved theorems / refuted claims / open formulations),
but localised at the level of formal terms within the formal register.  With
two readings of the operational register the middle and the refutable tier
are exhibited by the ZFC reading; the operational reading has only the
first tier and the trivial part of the third.
\end{remark}

%──────────────────────────────────────────────────────────────────────
\section{Methodological observations (Part IX)}
\label{sec:forms_method}

\subsection*{\S\,IX.1 Foundation-level properties}

\begin{remark}[Observation~1 --- The whole core module is import-light and axiom-free]
\label{rem:obs1_import}
\lean{VR.Forms.Register, VR.Forms.FormalTerm}
\leanok
\uses{def:Register, def:FormalTerm}
\texttt{Language.lean} introduces the syntactic skeleton without importing
any module at all.  Since 2026-09-12 the rest of the core module imports
only \texttt{VR.SetsOp} (itself Mathlib-free), and since 2026-09-13 the
whole core library imports nothing but Lean.  What was true of the base
layer in 2026-05 --- axiom-free \texttt{[]} --- is now true of every
declaration of VR-Forms in the core, by the build-time guard.
\end{remark}

\begin{remark}[Observation~2 --- Realisability inherits the axiom profile of its universe]
\label{rem:obs2_classical}
\lean{VR.Forms.isRealisable_empty, VR.Forms.isRealisable_omega,
      VR.Forms.isRealisable_pair}
\leanok
\uses{thm:isRealisable_base}
The realisable cases are exactly as clean as the universe that supplies
their witnesses.  Over \texttt{ZFSet} they were \texttt{[propext,
Quot.sound]} (the Classical-free closure layer of VR-Sets); over
\texttt{OpSet} they are \texttt{[]}.  Nothing in the realisability layer
itself introduces an axiom: it names witnesses and reads their defining
properties.
\end{remark}

\subsection*{\S\,IX.3 Structural patterns}

\begin{remark}[Observation~8 --- Universe handling across cross-cycle boundaries]
\label{rem:obs8_universe}
\lean{VR.Forms.isRealisable, VR.Forms.mixed_omega_two_register}
\leanok
\uses{def:isRealisable, thm:mixed_omega}
VR-Forms is a cross-cycle module, and universe inference does not
propagate across cycle boundaries without explicit annotation:
\texttt{OpSet.\{0\}} is pinned explicitly in \texttt{isRealisableN} and in
\texttt{Carrier} (as \texttt{OSet.\{0\}} was before), and the witnesses are
taken at that universe.
\end{remark}

\begin{remark}[Observation~9 --- The ZFA boundary is a boundary of groundedness]
\label{rem:obs9_zfa}
\lean{VR.Forms.mixed_AFA_two_registers}
\leanok
\uses{thm:mixed_AFA_two_registers}
\texttt{mixed\_AFA\_two\_registers} states in a single Lean \texttt{Prop}
what the two readings of the operational register say about one formal
description: realised in the operational universe, refuted in the
well-founded one.  What appeared in VR-Sets Part~X as Boundary~B.5 (``ZFA
total absence'') is thus not a limit of operationality but the choice of a
mode --- \texttt{IsGrounded}, a predicate on the graph, ZFC-mode inside
the one universe (Chapter~\ref{chap:setsop}).  The two-register apparatus is
precisely the instrument that lets this be said as a theorem rather than as
a remark.
\end{remark}

\subsection*{\S\,IX.4 Cross-cycle integration}

\begin{remark}[Observation~10 --- Two libraries, one perimeter]
\label{rem:obs10_zero_classical}
\lean{VR.Forms.bridge_AFA, VR.Forms.mixed_AFA_two_registers}
\leanok
The cycle plan predicted \texttt{[propext, Classical.choice, Quot.sound]}
or stricter; the 2026-05 result was \texttt{[propext, Quot.sound]} with zero
\texttt{Classical.choice}.  The 2026-09 result is sharper: every object of
VR-Forms that speaks about VR's own universe is in the core library
\texttt{VR} on \texttt{[]}, and the two theorems that speak about Mathlib's
universes (\texttt{bridge\_AFA}, the Conjecture bridges) together with the
junction that compares the two readings are in \texttt{VRClassical/Forms/}
\texttt{Bridge.lean}, at Mathlib's \texttt{[propext, Quot.sound]} and still
without \texttt{Classical.choice}.  The perimeter is drawn in code: nothing
in the core depends on the bridge.
\end{remark}

%──────────────────────────────────────────────────────────────────────
\section{Axiom profile}
\label{sec:forms_axiom_profile}

\axiomprofile{[]} for every public object of the core module
(\texttt{VR/Forms/}): the language, the classifier and both realisability
layers, all four base lemmas, the transit theorem, the three storeys of
conservativity, the four non-realisable terms, the positive mixed formula,
and the substrate results.  Enforced at build time by
\texttt{VR/Guard.lean}.

\axiomprofile{[propext, Quot.sound]} for the ZFC bridge
(\texttt{VRClassical/Forms/Bridge.lean}): \texttt{isRealisableZFC} and its
three positive cases, \texttt{bridge\_AFA}, the two Conjecture bridges,
\texttt{afa\_two\_registers} and \texttt{mixed\_AFA\_two\_registers}.
\textbf{Zero \texttt{Classical.choice}} anywhere in VR-Forms.

Comparison with VR-Sets (Chapter~\ref{chap:sets}): VR-Sets has 6 objects at
the full ceiling \texttt{[propext, Classical.choice, Quot.sound]} through
four structurally distinct mechanisms (ordinal-valued constructions,
definability, foundation, choice).  VR-Forms has none, and since the move
onto \texttt{OpSet} it has no axiom at all where it speaks for itself; the
only axioms left are those of the universe it is compared against.

%──────────────────────────────────────────────────────────────────────
\section{References}
\label{sec:forms_refs}

The preprint for this chapter is
\zenodo{10.5281/zenodo.20355939}{VR-Forms, v1.0.1} (23 May 2026).

The Lean~4 formalisation is archived at
\zenodo{10.5281/zenodo.20355757}{VR-Forms Lean, v1.0.0} (23 May 2026),
git tag \texttt{v1.3-vr-forms}; the move of the operational register onto
\texttt{OpSet} and the split into core and bridge are in the repository
history of 2026-09-12/13 (\texttt{CHANGELOG.md}).
